\documentclass[aspectratio=169,11pt]{beamer}
\usepackage[utf8]{inputenc}
\usepackage[T1]{fontenc}
\usepackage[spanish]{babel}
\usepackage{lmodern}
\usepackage{booktabs}
\usepackage{tabularx}
\usepackage{hyperref}
\usepackage{enumitem}

\usetheme{Madrid}
\usecolortheme{default}
\setbeamertemplate{navigation symbols}{}
\setbeamertemplate{footline}[frame number]

\title{IA agéntica para la creación de contenido académico}
\subtitle{Sesión 6 --- Presentaciones II: deck, notas y accesibilidad}
\author{Juliho Castillo}
\institute{CADi 40\,h · Escuela de Ingeniería y Ciencias\\
Tecnológico de Monterrey}
\date{Clave sugerida: \texttt{CADI-EIC-IA-AGENTICA-01}}

\begin{document}

\begin{frame}
\titlepage
\end{frame}

\begin{frame}{Agenda (3\,h) y producto E6}
\begin{tabularx}{\textwidth}{@{}lclX@{}}
\toprule
Bloque & Tiempo & Contenido \\
\midrule
Recap storyboards & 10 min & Priorizar 8--12 slides \\
Demo deck + notas & 25 min & Voz docente y variantes \\
Accesibilidad & 20 min & Contraste, ritmo, notación \\
Taller & 70 min & Deck MV + notas \\
Mini-revisión & 25 min & Checklist accesibilidad \\
Cierre & 20 min & Definir microartefacto final \\
\bottomrule
\end{tabularx}
\vspace{0.5em}
\textbf{E6:} deck MV + notas de orador + mini-checklist de accesibilidad.
\end{frame}

\begin{frame}{Resultados de aprendizaje (sesión)}
Al finalizar, usted podrá:
\begin{enumerate}
\item Construir un \textbf{deck MV} (8--12 diapositivas) desde el storyboard
\item Generar \textbf{notas de orador} alineadas a la voz docente
\item Aplicar una \textbf{mini-checklist} de accesibilidad
\item Declarar el \textbf{microartefacto final} EIC hacia el cierre del CADi
\end{enumerate}
\end{frame}

\begin{frame}{Del storyboard al deck MV}
\begin{itemize}
\item Materialice nodos: no invente slides ``porque quedan bonitas''
\item Máximo una idea fuerte por slide
\item Recorte si el tiempo real no alcanza
\item Versión MV: usable en clase, no perfecta
\end{itemize}
\begin{block}{CEDDIE}
Política de integridad / uso de IA: verificar por edición.
\end{block}
\end{frame}

\begin{frame}{Demo --- slides + notas de orador}
Mismo contenido, dos capas:
\begin{enumerate}
\item Slide: mensaje + visual/ecuación
\item Notas: 3--6 líneas de voz docente + tiempo + malentendido
\item Variante clase vs.\ conferencia (ritmo y profundidad)
\end{enumerate}
\begin{alertblock}{Mensaje clave}
Notas = continuidad de la voz docente, no párrafo pegado en la slide.
\end{alertblock}
\end{frame}

\begin{frame}{Variantes: clase vs.\ conferencia}
\begin{tabularx}{\textwidth}{@{}XX@{}}
\toprule
\textbf{Clase} & \textbf{Conferencia} \\
\midrule
Chequeos y pausas & Ritmo más continuo \\
Ejercicio breve & Caso ilustrativo \\
Más preguntas & Menos interrupciones \\
Pizarra / demo & Visuales autónomos \\
\bottomrule
\end{tabularx}
\vspace{0.8em}
Declare la variante en el brief del deck.
\end{frame}

\begin{frame}{Accesibilidad en decks técnicos}
Mini-checklist:
\begin{itemize}
\item[$\square$] Contraste texto/fondo suficiente
\item[$\square$] Tipografía y notación legibles a distancia
\item[$\square$] Orden de lectura predecible
\item[$\square$] Alt text conceptual para figuras clave
\item[$\square$] Ritmo: no bombardear en 20 s
\end{itemize}
Accesibilidad mejora aprendizaje para todos.
\end{frame}

\begin{frame}{Notación e inclusión}
\begin{itemize}
\item Evite color como único canal de significado
\item Explique símbolos en voz alta la primera vez
\item Figuras: diga qué debe verse (alt text conceptual)
\item Lenguaje claro; evite jerga innecesaria
\end{itemize}
\end{frame}

\begin{frame}{Anti-patrones de decks técnicos}
\begin{itemize}
\item Paredes de texto o derivaciones completas en una slide
\item Notas vacías (``hablar de la ecuación'')
\item Contraste pobre o símbolos ilegibles
\item 30 slides para 50 minutos
\item Decoración sin mensaje
\end{itemize}
\end{frame}

\begin{frame}{Integridad al publicar/impartir}
\begin{itemize}
\item Declarar asistencia de IA cuando corresponda
\item Verificar ecuaciones, unidades y citas antes de LMS
\item No pegar datos sensibles en prompts de generación
\item Política institucional $\rightarrow$ \textbf{CEDDIE} (vigencia por edición)
\end{itemize}
\end{frame}

\begin{frame}{Actividad de taller ($\sim$70 min)}
\textbf{Entregable E6:}
\begin{enumerate}
\item \textbf{Parte A --- Deck MV} (35 min): 8--12 slides desde E5
\item \textbf{Parte B --- Notas} (20 min): 3--6 líneas + tiempo + malentendido
\item \textbf{Parte C --- Accesibilidad} (15 min): mini-checklist completa
\end{enumerate}
Entrega: deck + \texttt{Apellido\_Sesion6\_notas-accesibilidad}
\end{frame}

\begin{frame}{Mini-checklist de accesibilidad (usar hoy)}
\begin{enumerate}
\item Contraste y tamaño
\item Notación legible
\item Orden de lectura
\item Alt text conceptual
\item Ritmo / tiempo por slide
\end{enumerate}
Marque N/A solo si justifica.
\end{frame}

\begin{frame}{Declaración de microartefacto final}
Antes del cierre del CADi (sesión 8), declare:
\begin{itemize}
\item Tipo: sección de libro / deck / artículo o guía
\item Audiencia EIC y alcance
\item Qué consolidará de sesiones 1--6
\item Pendientes de verificación
\end{itemize}
Hoy se fija el rumbo de acreditación.
\end{frame}

\begin{frame}{Rúbrica rápida (E6)}
\begin{tabularx}{\textwidth}{@{}lXX@{}}
\toprule
Criterio & Bien & A mejorar \\
\midrule
Fidelidad al storyboard & Nodos materializados & Slides inventadas \\
Notas & Voz docente usable & Vacías o muro \\
Accesibilidad & Checklist honesta & Ignorada \\
HITL & Edición visible & Generación ciega \\
\bottomrule
\end{tabularx}
\end{frame}

\begin{frame}{Takeaways de la sesión}
\begin{enumerate}
\item Deck MV = storyboard hecho realidad
\item Notas sostienen la voz docente
\item Accesibilidad $\neq$ trámite: es aprendizaje
\item Declare ya su microartefacto de sesión 8
\end{enumerate}
\end{frame}

\begin{frame}{Trabajo independiente}
Antes de la sesión 7:
\begin{itemize}
\item Avance microartefacto (bloque C)
\item Práctica de plantillas (bloque B)
\item Traer idea de \textbf{sección tipo artículo/guía} o consolidación
\end{itemize}
Parte de las \textbf{16\,h} independientes.
\end{frame}

\begin{frame}{Prep sesión 7 --- Artículos I}
Próxima sesión (3\,h):
\begin{itemize}
\item Estructura académica (IMRyD u otra)
\item Borrador asistido con \texttt{[VERIFICAR]}
\item Lista de citas/datos a verificar
\item Inicio del agent log del pipeline
\end{itemize}
\end{frame}

\begin{frame}{Gracias}
\begin{center}
{\Large Juliho Castillo}\\[0.5em]
Escuela de Ingeniería y Ciencias\\[1em]
CADi · IA agéntica para la creación de contenido académico\\[1.5em]
{\small \emph{Verificar políticas de integridad y uso de IA con CEDDIE antes de cada edición.}}
\end{center}
\end{frame}

\end{document}
